\section{Results}
\label{sec:results}

\subsection{Experimental setup}
\label{sec:results-setup}

We evaluate on six scRNA-seq datasets, each with a designated rare type: human immune dendritic cells (ASDC),
Baron pancreas (gamma), an integrated human pancreas (endothelial), and three Tabula Sapiens subsets --- lung
(lymphatic endothelial), small intestine (tuft) and stomach (mast). For every dataset we use a donor-held-out
$70/15/15$ split and vary the labelled rare fraction over $\{0.01,0.05,0.10,\text{all}\}$, repeating each
configuration over three random seeds ($42,43,44$). We compare scRareRefine against eight methods that read the
same split and the same labelled subset: the scANVI backbone, a $k$-NN classifier on the scANVI latent, CellTypist,
scBalance, ProtoCloud, scCAD, TOSICA, and HiCat. HiCat is transductive (its dimensionality reduction is fit on
combined train and test cells) and is therefore reported as an upper-bound reference rather than a fair inductive
competitor; TOSICA is run with a reduced pathway-token configuration for tractability. We report rare-cell F1 and the false-rescue
rate (FFR), aggregated as mean $\pm$ standard deviation over seeds. Because the labelled budget is
$\max(5,\lfloor p\,\Ntrain^{\rare}\rfloor)$, several nominal fractions collapse to the same labelled set on
small-rare datasets; we therefore report the label-scarce comparison over the $15$ \emph{distinct} labelled
configurations rather than the $18$ nominal ones.

\subsection{scRareRefine improves rare-cell F1 in the scarce regime with controlled FFR}
\label{sec:results-main}

\begin{figure*}[t]
\centering
\includegraphics[width=\textwidth]{fig2_recovery_curves.png}
\caption{\textbf{Label-scarcity recovery curves.} Rare-cell F1 for scANVI and scRareRefine across labelled rare-cell
fractions, shown as three-seed mean $\pm$ standard deviation for each dataset. The largest gains occur where scANVI
under-recovers the rare type at low label budgets, while scRareRefine converges to a no-op when the backbone is
already saturated.}
\label{fig:recovery}
\end{figure*}

\begin{figure*}[t]
\centering
\includegraphics[width=\textwidth]{fig2_main_summary.png}
\caption{\textbf{Main result.} (a) Mean rare-cell F1 in the label-scarce regime ($p\le0.10$), averaged over six
datasets, three fractions and three seeds. (b) Worst-case FFR over the whole benchmark; the dashed line is the
budget $\alpha=0.01$. scRareRefine is the only method that combines the highest scarce-regime F1 with a worst-case
FFR $\le\alpha$. $\dagger$ HiCat is transductive.}
\label{fig:main}
\end{figure*}

\begin{table}[t]
\centering
\small
\caption{\textbf{Per-dataset rare-cell F1 in the label-scarce regime} (mean over fractions $p\le0.10$ and three
seeds). scRareRefine is the best or tied-best inductive method on every dataset while keeping its worst-case FFR
within the budget $\alpha=0.01$. ``Best baseline'' excludes the transductive HiCat reference.}
\label{tab:main}
\begin{tabular}{lcccc}
\toprule
Dataset (rare type) & scANVI & Best baseline & \textbf{scRareRefine} & FFR$_{\max}$ \\
\midrule
immune dendritic (ASDC)       & 0.59 & 0.72 (kNN)        & \textbf{0.94} & 0.0010 \\
Baron pancreas (gamma)        & 0.53 & 0.46 (kNN)        & \textbf{0.66} & 0.0098 \\
lung (lymphatic EC)           & 0.72 & 0.92 (kNN)        & \textbf{0.98} & 0.0029 \\
small intestine (tuft)        & 0.98 & 0.98 (kNN)        & \textbf{0.98} & 0.0005 \\
stomach (mast)                & 0.61 & 0.49 (CellTypist) & \textbf{0.72} & 0.0002 \\
integrated pancreas (endo.)   & 0.89 & 0.98 (kNN)        & \textbf{1.00} & 0.0006 \\
\bottomrule
\end{tabular}
\end{table}

The recovery curves show the intended conditional behaviour before method-level aggregation (\cref{fig:recovery}).
At the most extreme budget, immune dendritic-cell F1 rises from $0.000$ to $0.927$, and the gain shrinks as more
rare labels are supplied ($+0.033$ at $p=0.10$). On Baron pancreas, scRareRefine improves the unstable scarce points
($0.383\!\to\!0.567$ at $p=0.01/0.05$; $0.820\!\to\!0.843$ at $p=0.10$), while on already-saturated integrated
pancreas and small intestine the gates mostly return the backbone. The headline is therefore label-scarcity recovery
rather than uniform improvement.

Aggregated across the label-scarce regime, scRareRefine attains the highest mean rare-cell F1 of all nine methods
(\cref{fig:main}a): $0.88$, against $0.72$ for the scANVI backbone and $0.73$ for the next-best inductive method
($k$-NN on the scANVI latent); the remaining methods range from $0.48$ to $0.65$, and the transductive HiCat
reference is lower ($0.19$) because it returns no rare cells on several datasets. Per dataset (\cref{tab:main}),
scRareRefine is the best or tied-best inductive method on all six datasets, with the largest gains where the backbone
is weakest (immune dendritic cells, $0.59\!\to\!0.94$; stomach mast cells, $0.61\!\to\!0.72$).

This F1 advantage is not bought with an uncontrolled false-positive cost. scRareRefine is the only method whose
worst-case FFR over the whole benchmark stays within budget ($0.0098\le\alpha$, \cref{fig:main}b); by contrast scCAD
($0.045$), scBalance ($0.016$) and TOSICA ($0.036$) exceed it. A paired analysis over the $54$ scarce
configurations (six datasets $\times$ three distinct fractions $\times$ three seeds) shows that scRareRefine never
underperforms the backbone: $29$ wins, $25$ ties --- the ties are abstentions --- and $0$ losses, with a mean F1
improvement of $+0.16$ (bootstrap $95\%$ CI $[+0.09,+0.24]$). Against every one of the eight baselines the
improvement's confidence interval excludes zero (one-sided Wilcoxon $p\le1.3\times10^{-6}$), which we read as
directional evidence given the correlated paired cells. Because the labelled budget has a floor of five cells, the
$18$ nominal scarce configurations collapse to $15$ distinct ones; over these, scRareRefine beats the majority of
baselines in $15/15$ and is the single best method in $14/15$ (Supplementary~S3).

\begin{figure*}[t]
\centering
\includegraphics[width=\textwidth]{fig6_comparison_heatmap.png}
\caption{\textbf{Per-method $\times$ per-dataset comparison in the label-scarce regime} ($p\le0.10$, mean over three
fractions and three seeds). (a) rare-cell F1 and (b) rare-cell recall, with $*$ marking the best method on each
dataset; (c) FFR, where darker shading means more false rescues. scRareRefine (boxed column) is best or tied-best in
F1 on five of six datasets while its FFR shading stays light throughout, whereas scCAD and TOSICA show the darkest
FFR cells. $\dagger$ HiCat is transductive (upper-bound reference).}
\label{fig:heatmap}
\end{figure*}

\begin{figure*}[t]
\centering
\includegraphics[width=\textwidth]{fig7_comparison_boxplot.png}
\caption{\textbf{Across-dataset distribution in the label-scarce regime.} Each box summarises a method over $54$
scarce configurations (six datasets $\times$ three fractions $\times$ three seeds); the median is labelled. Top to
bottom: rare-cell F1, rare-cell recall, FFR (dashed line is the budget $\alpha=0.01$). scRareRefine has the highest
and tightest F1/recall distributions and is the only method whose FFR box stays at zero, below the budget. $\dagger$
HiCat is transductive.}
\label{fig:box}
\end{figure*}

\Cref{fig:heatmap} resolves the per-method summary of \cref{fig:main} into the full method$\times$dataset grid:
scRareRefine carries the F1/recall stars on most datasets while keeping uniformly light FFR shading, in contrast to
the dark FFR cells of scCAD and TOSICA. \Cref{fig:box} shows the same effect as across-dataset distributions ---
scRareRefine attains the highest and tightest F1 and recall medians ($0.93$ and $0.93$) and is the only method whose
FFR box rests at zero, below $\alpha$.

\subsection{Component ablation}
\label{sec:results-ablation}

\begin{figure*}[t]
\centering
\includegraphics[width=0.49\textwidth]{fig3a_ablation_components.png}\hfill
\includegraphics[width=0.49\textwidth]{fig3b_ablation_rank.png}
\caption{\textbf{Ablation.} Left: leave-one-out over the four components (F1 top, worst-case FFR bottom). Right:
candidate-rank sensitivity (fixed $k$ vs.\ validation-adaptive). Removing the conformal threshold or the
separability gate barely changes mean F1 but breaks the FFR budget; the adaptive rank is the component that raises
F1; a fixed rank $=3$ over-fires FFR.}
\label{fig:ablation}
\end{figure*}

The full method raises mean rare-cell F1 from $0.76$ (backbone) to $0.89$ (\cref{fig:ablation}, left). Removing each
component in turn separates two roles. The validation-adaptive candidate-rank selector is the only component that
materially raises F1: replacing it with the most conservative fixed rank ($k=1$) lowers F1 by $0.010$. The conformal
threshold and the two gates, by contrast, barely change mean F1 but are what keep the method safe. Removing the
conformal threshold leaves mean F1 essentially unchanged ($+0.002$) yet pushes the worst-case FFR to $0.0165$, above
budget; removing the separability gate even raises mean F1 (by $0.018$) but likewise breaks the budget (worst-case
FFR $0.0153$); and removing the necessity guard leaves the benchmark mean unchanged but causes per-dataset
regressions on already-saturated datasets (integrated pancreas, $0.990\!\to\!0.977$, with a non-zero FFR appearing).
The rank-sensitivity panel (\cref{fig:ablation}, right) confirms the value of adapting $k$: every fixed choice is
dominated in F1 by the adaptive selector ($k{=}1$: $0.877$, $k{=}2$: $0.865$, $k{=}3$: $0.853$ vs.\ adaptive
$0.887$), and the most permissive fixed rank ($k{=}3$) drives the worst-case FFR to $0.046$, far above budget,
whereas the adaptive selector stays within it ($0.0098$). In summary, the adaptive rank drives F1, the conformal
threshold controls FFR, and the two gates prevent regressions.

\subsection{The separability gate is a conservative, dataset-dependent prior}
\label{sec:results-sep}

\begin{figure*}[t]
\centering
\includegraphics[width=\textwidth]{fig4_separability.png}
\caption{\textbf{Separability-gate analysis.} (a,b) Controlled semi-synthetic stress test (entangling a rare type
toward its nearest majority): rescue stays FFR-safe well below the gate at $1.3$. (c) Benchmark-wide sensitivity of
the gate threshold: $1.3$ is the smallest threshold that keeps worst-case FFR $\le\alpha$ on real data.}
\label{fig:sep}
\end{figure*}

The separability gate uses a fixed, dataset-independent threshold ($\sep_{\mathrm{low}}=1.3$). We probe this choice
from two directions (\cref{fig:sep}). First, a controlled semi-synthetic stress test progressively entangles a rare
population (lung lymphatic endothelial cells) toward its nearest majority type; rescue \emph{without} the gate stays
within the FFR budget down to separability $\approx0.76$ and produces a single marginal violation ($18$ false
rescues out of $1716$, FFR $0.0105$) only at the lowest separability ($\approx0.69$), well below the gate
(\cref{fig:sep}a,b). Second, a benchmark-wide sweep of the gate threshold (\cref{fig:sep}c) shows that $1.3$ is the
smallest threshold that keeps the worst-case FFR within budget on real data: lowering it to $\le1.0$ raises mean F1
by $0.018$ but admits two FFR violations (worst-case $0.0153$, on the low-separability Baron-pancreas configurations
at $\sep\approx1.22$), whereas raising it to $1.6$ only sacrifices F1 ($0.888\!\to\!0.857$) with no further FFR
benefit. The two experiments disagree on where rescue becomes unsafe ($\approx1.22$ on real Baron pancreas vs.\
$\approx0.69$ on the synthetic sweep), which is the central observation: separability alone does not monotonically
determine rescue risk. We therefore present $1.3$ not as a located collapse boundary but as a deliberately
conservative, pre-specified prior chosen by the FFR constraint rather than tuned for F1 --- lowering it would in fact
\emph{raise} mean F1 --- that occasionally abstains where rescue would have been safe.

\subsection{Qualitative rescue in latent space}
\label{sec:results-umap}

\begin{figure*}[t]
\centering
\includegraphics[width=\textwidth]{fig5_umap_immune.png}
\caption{\textbf{Rescue in the scANVI latent space (immune dendritic cells).} Cells missed by scANVI that lie near
the rare cluster are recovered (recall $0.72\to0.90$), whereas geometrically distant rare cells remain missed,
illustrating the recall ceiling discussed in \cref{sec:discussion}.}
\label{fig:umap}
\end{figure*}

\cref{fig:umap} visualises the rescue on immune dendritic cells in the scANVI latent space. Of the true rare cells
(\cref{fig:umap}a), the backbone recovers those forming a compact cluster (recall $0.72$, panel b). scRareRefine
relabels $23$ further cells lying at the edge of this cluster, raising recall to $0.90$ at precision $0.99$, with a
single false rescue (panels c,d). The $13$ rare cells that remain missed are scattered far from the rare prototype;
admitting them would require widening the candidate pool past the FFR budget. This is the geometric recall ceiling
discussed in \cref{sec:discussion}.
